\documentclass[11pt]{article}
\usepackage[margin=1in]{geometry}
\usepackage{amsmath,amssymb,amsthm,mathtools}
\usepackage{microtype}
\usepackage[hidelinks]{hyperref}
\usepackage{enumitem}

\newtheorem{theorem}{Theorem}[section]
\newtheorem{lemma}[theorem]{Lemma}
\newtheorem{corollary}[theorem]{Corollary}
\newtheorem{proposition}[theorem]{Proposition}
\theoremstyle{remark}
\newtheorem{remark}[theorem]{Remark}

\newcommand{\E}{\mathbb E}
\newcommand{\Pp}{\mathbb P}
\newcommand{\Var}{\operatorname{Var}}
\newcommand{\Cov}{\operatorname{Cov}}
\newcommand{\1}{\mathbf 1}
\newcommand{\N}{\mathbb N}
\newcommand{\R}{\mathbb R}

\title{Degree Fluctuations in Bilateral Preference Random Graphs}
\author{Yitzchak Shmalo}
\date{September 19, 2026}

\begin{document}
\maketitle

\begin{abstract}
Let $LK(n,k)$ be the bilateral preference random graph obtained by placing independent continuous scores on the edges of $K_n$, letting each vertex prefer its $k$ incident edges of smallest score, and retaining an edge precisely when both endpoints prefer it. Dabirian and Subramanian asked for the limiting distribution of the standardized degree $(D-\E D)/\sqrt{\Var D}$ and observed that even $\Var D$ is difficult because of the dependence between edge-retention events. We characterize the degree fluctuations in all asymptotic regimes. If $k\to\infty$ and $n-k\to\infty$, then, with
\[
s_n^2=\frac{k(n-k)}n,
\qquad g(z)=\phi(z)+z\Phi(z),
\]
we prove
\[
\frac{k-D}{s_n}\Longrightarrow g(Z)
\]
with convergence of the first two moments. Hence
\[
\E D=k-\frac{s_n}{\sqrt\pi}+o(s_n),
\qquad
\Var(D)=\left(\frac13+\frac{\sqrt3-2}{2\pi}\right)s_n^2+o(s_n^2),
\]
and the standardized degree converges to an explicit non-Gaussian law. At the opposite dense boundary, if $r=n-1-k\ge1$ is fixed, then $k-D$ converges to an explicit Gamma-mixed Poisson variable. These boundary laws vary with $r$, while as $r\to\infty$ they converge after standardization to the same interior non-Gaussian law. Fixed finite $k$ is reconciled with the previously known Erlang-weighted-tree local limit. Thus there is a universal interior law but no single standardized law for arbitrary switching sequences $k(n)$.
\end{abstract}

\section{Model and source question}

Assign independent $U_{ij}=U_{ji}\sim\mathrm{Unif}(0,1)$ to the edges of $K_n$. Vertex $i$ prefers the $k$ incident edges of smallest score. The graph $LK(n,k)$ retains $(i,j)$ if and only if both $i$ and $j$ prefer that edge. Since only score ranks matter, this is equivalent to the score orientation used in the source literature.

Fix a root vertex $0$ and write $D=D_{n,k}$ for its degree. Dabirian and Subramanian ask for the limiting distribution of
\[
\frac{D-\E D}{\sqrt{\Var D}},
\]
and note that even computing $\Var D$ is difficult because the bilateral decisions are dependent. Their sharp mean-degree theorem is stated in the growing sparse regime $k\to\infty$, $k=o(n^{1/3})$. Our first theorem strictly contains that regime and, in fact, covers every interior sequence $k\to\infty$, $n-k\to\infty$.

Put
\[
H:=k-D,
\qquad
s_n^2:=\frac{k(n-k)}n,
\qquad
p_n:=\frac kn,
\qquad
\sigma_n:=\frac{s_n}{n}.
\]
Thus $H$ counts root-preferred edges rejected by their opposite endpoints.

Let $\phi$ and $\Phi$ denote the standard normal density and distribution function, and define
\[
g(z):=\phi(z)+z\Phi(z)=\E[(z-Z')_+],
\]
where $Z'\sim N(0,1)$. Finally let
\[
v:=\Var(g(Z))=\frac13+\frac{\sqrt3-2}{2\pi}.
\]

\section{Interior growing regime}

\begin{theorem}[Universal interior law]\label{thm:interior}
Suppose
\[
k=k(n)\to\infty,
\qquad
n-k\to\infty.
\]
Then
\[
\frac{H}{s_n}\Longrightarrow g(Z)
\]
and the first two moments converge. Consequently,
\[
\E D=k-\frac{s_n}{\sqrt\pi}+o(s_n),
\]
\[
\Var(D)=v s_n^2+o(s_n^2),
\]
and
\[
\frac{D-\E D}{\sqrt{\Var D}}
\Longrightarrow
L_\infty:=\frac{\pi^{-1/2}-g(Z)}{\sqrt v}.
\]
\end{theorem}

\subsection{Exact order-statistic representation}

Let $S$ be the set of the root's $k$ preferred neighbors and let $X$ be the $(k+1)$st smallest among the $n-1$ root-edge scores. Conditional on $(X,S)$, the selected root-edge scores $(U_i)_{i\in S}$ are i.i.d. $\mathrm{Unif}(0,X)$, and all non-root edge scores remain independent uniforms.

For $i\in S$, let $T_i$ be the $k$th smallest among the $n-2$ scores incident to $i$ other than $(0,i)$. Then
\[
T_i\sim \mathrm{Beta}(k,n-k-1),
\]
and the root edge $(0,i)$ is rejected by $i$ exactly when $U_i\ge T_i$. If $T$ is an independent $\mathrm{Beta}(k,n-k-1)$ variable, define
\[
q_n(x):=\frac1x\E[(x-T)_+].
\]
With $R_i:=\1_{\{U_i\ge T_i\}}$ we therefore have the exact identity
\[
\E[H\mid X]=kq_n(X).
\tag{2.1}
\]

\subsection{The dependence localizes to one shared edge}

The key point is that the dependence between two indicators $R_i$ and $R_j$ is one-dimensional.

\begin{lemma}[One-edge covariance bound]\label{lem:cov}
For distinct $i,j\in S$,
\[
\Cov(R_i,R_j\mid X)
\le
\frac{1}{4(n-2)^2X^2}.
\]
\end{lemma}

\begin{proof}
The non-root stars of $i$ and $j$ share exactly one edge, $(i,j)$; call its score $W$. Conditional on $(X,W)$, the two indicators depend on disjoint remaining edge sets and independent root-edge scores, so they are conditionally independent. Hence
\[
\Cov(R_i,R_j\mid X)=\Var(p_X(W)\mid X),
\]
where $p_X(w)=\Pp(R_i=1\mid X,W=w)$.

Condition further on the root-edge score $U_i=u$. After deleting $W$, let
\[
B_u\sim\mathrm{Bin}(n-3,u)
\]
count the remaining competitors below $u$. Rejection means
\[
B_u+\1_{\{w<u\}}\ge k.
\]
Thus, for $0\le w<X$,
\[
p_X(w)=b_X+\frac1X\int_w^X
\Pp(\mathrm{Bin}(n-3,u)=k-1)\,du,
\]
where $b_X$ does not depend on $w$. Therefore
\[
\operatorname{osc}(p_X)
\le \frac1X\int_0^1
\binom{n-3}{k-1}u^{k-1}(1-u)^{n-k-2}\,du.
\]
The beta integral is exactly $1/(n-2)$, hence
\[
\operatorname{osc}(p_X)\le \frac{1}{(n-2)X}.
\]
Popoviciu's inequality gives the claimed conditional variance bound.
\end{proof}

Since $\Var(R_i\mid X)\le q_n(X)$, Lemma~\ref{lem:cov} yields
\[
\Var(H\mid X)
\le kq_n(X)+\frac{k(k-1)}{4(n-2)^2X^2}.
\tag{2.2}
\]

Let
\[
G_n:=\left\{X\ge \frac{k}{2n}\right\}.
\]
Since $X$ is the $(k+1)$st order statistic among $n-1$ uniforms,
\[
\Pp(G_n^c)
=
\Pp\left(\mathrm{Bin}\left(n-1,\frac{k}{2n}\right)\ge k+1\right)
\le e^{-c k}
\tag{2.3}
\]
for the universal constant $c=\log 2-1/2>0$.

On $G_n$ the covariance contribution in (2.2) is $O(1)$. Moreover,
\[
kq_n(X)
\le 2n\E(|X-T|\mid X),
\]
so, by Cauchy--Schwarz and the beta variance formulas,
\[
\E[kq_n(X);G_n]=O(s_n+1).
\]
On $G_n^c$ we use $0\le H\le k$ and the exponential bound (2.3). Since $s_n\to\infty$, we obtain
\[
\E\Var(H\mid X)=o(s_n^2).
\]
Equivalently,
\[
\frac{H-kq_n(X)}{s_n}\longrightarrow0
\quad\text{in }L^2.
\tag{2.4}
\]

\subsection{Asymptotics of the conditional mean}

Set
\[
U_n:=\frac{X-p_n}{\sigma_n},
\qquad
V_n:=\frac{T-p_n}{\sigma_n}.
\]
Here
\[
X\sim\mathrm{Beta}(k+1,n-k-1),
\qquad
T\sim\mathrm{Beta}(k,n-k-1),
\]
independently. Because both beta shape parameters tend to infinity,
\[
U_n\Longrightarrow Z,
\qquad
V_n\Longrightarrow Z',
\]
with independent standard-normal limits. Direct beta moment formulas give uniformly bounded fourth moments.

Define
\[
g_n(u):=\E[(u-V_n)_+].
\]
The convergence $V_n\Rightarrow Z'$ plus uniform integrability implies convergence in Wasserstein distance $W_1$, and since $x\mapsto(u-x)_+$ is $1$-Lipschitz uniformly in $u$,
\[
\sup_{u\in\R}|g_n(u)-g(u)|\to0.
\tag{2.5}
\]
The exact conditional mean rescales as
\[
\frac{kq_n(X)}{s_n}
=
\frac{k}{nX}\,g_n(U_n).
\tag{2.6}
\]
Also $k/(nX)=p_n/X\to1$ in probability. Combining (2.5)--(2.6),
\[
\frac{kq_n(X)}{s_n}\Longrightarrow g(Z).
\tag{2.7}
\]
The good-event estimate $k/(nX)\le2$ on $G_n$, the fourth-moment bounds for $U_n,V_n$, and the exponentially small complement $G_n^c$ give uniform integrability of the square of the left side of (2.7). Thus the first two moments converge.

Combining (2.4) and (2.7) proves
\[
H/s_n\Rightarrow g(Z)
\]
with first-two-moment convergence.

It remains only to identify the constants. With $Z,Z_1,Z_2$ i.i.d. standard normal,
\[
\E g(Z)=\E[(Z-Z_1)_+]=\frac1{\sqrt\pi},
\]
and
\[
\E g(Z)^2
=
\E[(Z-Z_1)_+(Z-Z_2)_+]
=
\frac13+\frac{\sqrt3}{2\pi}.
\]
Hence
\[
v=\frac13+\frac{\sqrt3-2}{2\pi}.
\]
Theorem~\ref{thm:interior} follows.

\begin{corollary}[Direct answer in the source's mean-theorem regime]\label{cor:source}
If $k\to\infty$ and $k=o(n^{1/3})$, then
\[
\frac{D-\E D}{\sqrt{\Var D}}\Longrightarrow L_\infty.
\]
In particular, the limiting standardized degree in the regime surrounding the source's Theorem 8 is non-Gaussian.
\end{corollary}

\begin{proof}
The hypotheses imply $n-k\to\infty$, so Theorem~\ref{thm:interior} applies directly.
\end{proof}

\section{An explicit form of the interior limiting law}

Since
\[
g'(z)=\Phi(z)>0,
\]
$g$ is strictly increasing from $(0,\infty)$ as $z$ runs from $-\infty$ to $\infty$. Put
\[
m:=\pi^{-1/2}.
\]
The support of
\[
L_\infty=\frac{m-g(Z)}{\sqrt v}
\]
is $(-\infty,m/\sqrt v)$. For $x<m/\sqrt v$, let
\[
y_x:=g^{-1}(m-\sqrt v\,x).
\]
Then
\[
\Pp(L_\infty\le x)=1-\Phi(y_x),
\tag{3.1}
\]
while the distribution function equals $1$ for $x\ge m/\sqrt v$. Differentiating (3.1) gives the density
\[
f_\infty(x)
=
\frac{\sqrt v\,\phi(y_x)}{\Phi(y_x)},
\qquad x<m/\sqrt v.
\tag{3.2}
\]
This gives a completely explicit one-dimensional characterization of the source's limiting law: the only inversion required is that of the smooth strictly increasing function $g$.

The bounded upper support already rules out Gaussianity.

\section{Fixed-complement boundary}

We now let
\[
r:=n-1-k
\]
remain fixed. Thus every vertex rejects exactly its $r$ largest incident scores.

For the root-edge $(0,i)$ define the upper-tail coordinate
\[
Y_i^{(n)}:=n(1-U_{0i}).
\]
Let $A_{n,r}$ be the $r$th smallest among the $n-1$ coordinates $Y_i^{(n)}$. The root proposes $(0,i)$ exactly when
\[
Y_i^{(n)}>A_{n,r}.
\]
Let $J_i^{(n)}$ indicate that vertex $i$ rejects the root edge. Conditional on $Y_i^{(n)}=y$,
\[
\Pp(J_i^{(n)}=1\mid Y_i^{(n)}=y)
=b_{n,r}(y)
:=
\Pp(\mathrm{Bin}(n-2,y/n)\le r-1).
\tag{4.1}
\]
Therefore
\[
H_n
=
\sum_{i\ne0}
\1_{\{Y_i^{(n)}>A_{n,r}\}}J_i^{(n)}.
\tag{4.2}
\]

\subsection{Marked Poisson convergence}

The point process
\[
\Pi_n:=\sum_{i\ne0}\delta_{Y_i^{(n)}}
\]
converges to a rate-one Poisson point process $\Pi$ on $(0,\infty)$. Hence
\[
A_{n,r}\Longrightarrow A_r,
\qquad
A_r\sim\Gamma(r,1).
\tag{4.3}
\]
For fixed $y$,
\[
b_{n,r}(y)\to p_r(y)
:=
 e^{-y}\sum_{j=0}^{r-1}\frac{y^j}{j!}
=\Pp(\Gamma(r,1)>y).
\tag{4.4}
\]

The marks are not exactly independent because two candidate vertices share one edge. The next lemma shows that this dependence disappears on every bounded upper-tail window.

\begin{lemma}[Asymptotic independent marking]\label{lem:marking}
Fix $L<\infty$ and distinct labels $i_1,\ldots,i_m$. Conditional on
\[
Y_{i_a}^{(n)}=y_a\in[0,L],\qquad 1\le a\le m,
\]
the mark vector $(J_{i_1}^{(n)},\ldots,J_{i_m}^{(n)})$ converges uniformly on $[0,L]^m$ to independent Bernoulli variables with success probabilities $p_r(y_1),\ldots,p_r(y_m)$.
\end{lemma}

\begin{proof}
Delete the $\binom m2$ edges joining two candidate vertices. For candidate $i_a$, the number of remaining non-root incident edges larger than its root-edge score is
\[
C_a^{(n)}\sim\mathrm{Bin}(n-m-1,y_a/n),
\]
and these variables are independent across $a$. Uniformly for $y_a\in[0,L]$, they converge to independent $\mathrm{Poisson}(y_a)$ variables.

A deleted candidate--candidate edge can change a mark only if its score exceeds at least one of the relevant root-edge thresholds. For one such internal edge that probability is at most $L/n$. A union bound therefore gives an error at most
\[
\binom m2\frac Ln=o(1).
\]
On the complementary event, candidate $i_a$ rejects the root edge exactly when $C_a^{(n)}\le r-1$. This proves the uniform factorization and the stated Bernoulli limits.
\end{proof}

To remove the compact cutoff, note that for fixed $r$ and all sufficiently large $n$,
\[
b_{n,r}(y)
\le
 e^{-y/2}\sum_{j=0}^{r-1}\frac{y^j}{j!},
\qquad 0\le y\le n.
\tag{4.5}
\]
The right side is integrable on $[0,\infty)$. Thus marked points beyond $L$ contribute an expected amount tending uniformly to zero as $L\to\infty$. The same envelope, applied to two marks after conditioning on their unique shared edge, controls the second factorial moment as well. Hence distribution, mean and variance all converge.

\begin{theorem}[Fixed-complement mixed-Poisson law]\label{thm:fixedr}
Fix an integer $r\ge1$ and put $k=n-1-r$. Then
\[
H_n\Longrightarrow H_r
\]
with convergence of the first two moments, where
\[
A_r\sim\Gamma(r,1),
\qquad
H_r\mid A_r=a
\sim\mathrm{Poisson}(\Lambda_r(a)),
\]
and
\[
\Lambda_r(a)
=
\int_a^\infty p_r(y)\,dy
=
 e^{-a}\sum_{\ell=0}^{r-1}(r-\ell)\frac{a^\ell}{\ell!}.
\tag{4.6}
\]
Equivalently,
\[
\E[s^{H_r}]
=
\frac1{(r-1)!}
\int_0^\infty
 a^{r-1}e^{-a}
\exp\!\bigl(-(1-s)\Lambda_r(a)\bigr)\,da,
\qquad 0\le s\le1.
\tag{4.7}
\]
\end{theorem}

\begin{proof}
By Lemma~\ref{lem:marking}, on compact windows the limiting root point process is independently marked with location-dependent mark probability $p_r(y)$. Conditional on $A_r=a$, the points of $\Pi$ after $a$ form an independent rate-one Poisson process on $(a,\infty)$. Independent thinning gives a Poisson number of accepted marks with mean $\int_a^\infty p_r(y)\,dy$. The incomplete-gamma identity gives (4.6). The integrable envelope (4.5), together with the corresponding two-point bound, removes the truncation and yields first-two-moment convergence. Formula (4.7) is the probability generating function of the resulting mixed-Poisson variable.
\end{proof}

\subsection{Exact moments and an independent beta--gamma derivation}

Let
\[
\mu_r:=\E H_r,
\qquad
\nu_r:=\Var(H_r).
\]
Since $H_r\mid A_r$ is Poisson,
\[
\mu_r=\E\Lambda_r(A_r),
\qquad
\nu_r=\mu_r+\Var(\Lambda_r(A_r)).
\tag{4.8}
\]
A direct finite-sum evaluation gives
\[
\mu_r
=r\frac{\binom{2r}{r}}{4^r}
=
\frac{\Gamma(r+1/2)}{\sqrt\pi\,\Gamma(r)}.
\tag{4.9}
\]
There is also a useful independent derivation. Let $A_r,G_r$ be i.i.d. $\Gamma(r,1)$. Then
\[
\Lambda_r(A_r)=\E[(G_r-A_r)_+\mid A_r],
\]
so
\[
\mu_r=\E[(G_r-A_r)_+]=\frac12\E|G_r-A_r|.
\tag{4.10}
\]
Write
\[
S=A_r+G_r,
\qquad
B=\frac{A_r}{A_r+G_r}.
\]
By beta--gamma algebra, $S\sim\Gamma(2r,1)$ and $B\sim\mathrm{Beta}(r,r)$ are independent. Since
\[
(G_r-A_r)_+=S(1-2B)_+,
\]
we get
\[
\mu_r
=2r\,\E[(1-2B)_+]
=2r\frac{1}{B(r,r)}\int_0^{1/2}(1-2b)b^{r-1}(1-b)^{r-1}\,db,
\]
which evaluates to (4.9). This provides a second check on the moment formula without using the polynomial expression (4.6).

For the variance, put
\[
M_{2,r}:=\E[\Lambda_r(A_r)^2].
\]
Expanding (4.6) and integrating against the $\Gamma(r,1)$ density gives the exact finite sum
\[
M_{2,r}
=
\frac1{(r-1)!}
\sum_{\ell,m=0}^{r-1}
\frac{(r-\ell)(r-m)(r+\ell+m-1)!}
{\ell!\,m!\,3^{r+\ell+m}}.
\tag{4.11}
\]
Hence
\[
\nu_r=\mu_r+M_{2,r}-\mu_r^2>0.
\tag{4.12}
\]
Moment convergence in Theorem~\ref{thm:fixedr} yields
\[
\E D=k-\mu_r+o(1),
\qquad
\Var(D)=\nu_r+o(1),
\]
and therefore
\[
\frac{D-\E D}{\sqrt{\Var D}}
\Longrightarrow
L_r:=\frac{\mu_r-H_r}{\sqrt{\nu_r}}.
\tag{4.13}
\]

For example,
\[
r=1:\qquad \mu_1=\frac12,\qquad \nu_1=\frac7{12},
\]
while
\[
r=2:\qquad \mu_2=\frac34,\qquad \nu_2=\frac{433}{432}.
\]
Since $\Pp(H_r=0)>0$, the upper endpoint of $L_r$ is $\mu_r/\sqrt{\nu_r}$. Its square is $3/7$ for $r=1$ and $243/433$ for $r=2$, so $L_1$ and $L_2$ are distinct.

\section{Boundary-to-interior gluing}

The finite-complement family joins continuously to the interior law.

\begin{theorem}[Gluing]\label{thm:glue}
As $r\to\infty$,
\[
\frac{H_r}{\sqrt r}\Longrightarrow g(Z)
\]
with convergence of the first two moments. Consequently,
\[
\frac{\mu_r}{\sqrt r}\to\frac1{\sqrt\pi},
\qquad
\frac{\nu_r}{r}\to v,
\qquad
L_r\Longrightarrow L_\infty.
\]
\end{theorem}

\begin{proof}
Let $A_r,G_r$ be i.i.d. $\Gamma(r,1)$. From (4.10), conditionally on $A_r$,
\[
\Lambda_r(A_r)=\E[(G_r-A_r)_+\mid A_r].
\]
Set
\[
U_r:=\frac{A_r-r}{\sqrt r},
\qquad
V_r:=\frac{G_r-r}{\sqrt r}.
\]
The gamma CLT gives $U_r,V_r\Rightarrow N(0,1)$, and their second moments are uniformly bounded. Hence the law of $V_r$ converges to $N(0,1)$ in $W_1$. Therefore
\[
\frac{\Lambda_r(A_r)}{\sqrt r}
=
\E[(V_r-U_r)_+\mid U_r]
\Longrightarrow g(Z)
\]
in distribution, after replacing $Z$ by $-Z$ if desired; the two have the same law.

The conditional Poisson noise is negligible on the $\sqrt r$ scale:
\[
\frac1r\E[(H_r-\Lambda_r(A_r))^2]
=
\frac{\mu_r}{r}\to0,
\]
where $\mu_r\asymp\sqrt r$ by (4.9). Thus $H_r/\sqrt r\Rightarrow g(Z)$. The same estimates and the exact moment formulas give first-two-moment convergence, so standardization yields $L_r\Rightarrow L_\infty$.
\end{proof}

\section{Regime-complete interpretation}

The remaining asymptotic regime is fixed finite $k$. Moharrami, Subramanian, Liu and Sundaresan prove local weak convergence of the finite bilateral-preference model, with fixed potential degree distribution, to the Erlang weighted tree and explicitly identify the root-degree law. Specializing the potential degree distribution to a point mass at a fixed integer $k=\kappa$ gives the fixed-$\kappa$ degree limit relevant here. Since $0\le D\le\kappa$, local weak convergence automatically gives convergence of the first two moments as well.

We can therefore summarize all subsequential degree limits.

\begin{theorem}[Regime classification]\label{thm:classification}
Let $k=k(n)$ with $1\le k\le n-1$, and consider any subsequence along which the standardized degree is defined. After passing to a further subsequence, exactly one of the following descriptions applies:
\begin{enumerate}[label=(\roman*)]
\item $k\equiv\kappa<\infty$: the standardized limit is the standardized root-degree law of the fixed-$\kappa$ Erlang weighted tree;
\item $k\to\infty$ and $n-k\to\infty$: the standardized limit is $L_\infty$ from Theorem~\ref{thm:interior};
\item $r=n-1-k\equiv r_0\ge1$: the standardized limit is $L_{r_0}$ from Theorem~\ref{thm:fixedr};
\item $k=n-1$: the graph is complete, $D=n-1$ deterministically, and standardization is undefined because $\Var D=0$.
\end{enumerate}
In particular, there is no single standardized limiting law valid for arbitrary switching sequences $k(n)$.
\end{theorem}

\begin{proof}
Every integer sequence $k(n)$ admits a further subsequence on which either $k$ is bounded or $k\to\infty$. In the bounded case, a further subsequence makes $k$ constant, giving (i). In the growing case, either $n-k\to\infty$ along a further subsequence, giving (ii), or $n-k$ is bounded. In the latter case, since $n-1-k$ is integer and nonnegative, a further subsequence makes it constant, giving (iii) or (iv).

To see that a universal law cannot exist without assumptions on $k(n)$, alternate between complements $r=1$ and $r=2$. The standardized subsequential limits are $L_1$ and $L_2$, which are distinct because their upper support endpoints differ.
\end{proof}

\section{Why the dependence is tractable}

The source correctly identifies dependence as the main obstacle. The proofs show that there are two distinct mechanisms, each local after the right conditioning.

In the interior regime, conditioning on the root cutoff $X$ leaves the root-retention indicators dependent only through one shared non-root edge per pair. The exact beta integral in Lemma~\ref{lem:cov} makes the total pair-covariance contribution $O(1)$ on the high-probability event $X\asymp k/n$. Thus all conditional edge noise is lower order than the random cutoff fluctuation.

At fixed complement, the correct scale is not Gaussian around $k/n$ but the endpoint point process $n(1-U)$. On every bounded window, the finitely many candidate vertices share only finitely many internal edges, each of which matters with probability $O(1/n)$. The dependence therefore vanishes in the marked-Poisson limit.

These are the same structural statement at two different scales: bilateral dependence is global in appearance but pairwise-local after conditioning on the root's order statistic.

\section{Prior-art scope}

Dabirian and Subramanian prove connectivity results and a sharp asymptotic expansion for the mean degree, and explicitly leave the standardized degree distribution as an open question. Moharrami, Subramanian, Liu and Sundaresan analyze the fixed-potential-degree local weak limit through the Erlang weighted tree and characterize its root-degree law; that fixed-$k$ result is therefore treated here as prior art, not as a new theorem.

A bounded literature check dated September 19, 2026 found the final Cambridge version of the source, the Erlang weighted tree paper, the original La--Kabkab model paper, and no later indexed paper giving either the interior law of Theorem~\ref{thm:interior} or the fixed-complement mixed-Poisson law of Theorem~\ref{thm:fixedr}. Cambridge currently reports zero Crossref citations for the source article. This is only bounded negative evidence and is not a novelty certificate.

\section{Conclusion}

The standardized degree question has a clean but non-Gaussian answer. Throughout the entire interior growing regime, the deficit $k-D$ is asymptotically the positive-part Gaussian smoothing $g(Z)$ at scale $\sqrt{k(n-k)/n}$. At finite complement the limit becomes an explicit mixed-Poisson variable, and those boundary laws converge back to the same interior distribution as the complement grows. Together with known fixed-$k$ local weak limits, this yields a regime-complete description and explains precisely how the complicated bilateral dependence disappears after the appropriate order-statistic conditioning.

\begin{thebibliography}{9}

\bibitem{DS}
H. Dabirian and V. Subramanian,
\emph{Connectivity of a family of bilateral preference random graphs},
Journal of Applied Probability 63 (2026), no.~2, 552--579.
DOI: \href{https://doi.org/10.1017/jpr.2025.10051}{10.1017/jpr.2025.10051}.

\bibitem{EWT}
M. Moharrami, V. Subramanian, M. Liu, and R. Sundaresan,
\emph{The Erlang weighted tree, a new branching process},
Random Structures \& Algorithms 64 (2024), 537--624.
DOI: \href{https://doi.org/10.1002/rsa.21180}{10.1002/rsa.21180}.

\bibitem{LK}
R. J. La and M. Kabkab,
\emph{A new random graph model with self-optimizing nodes: Connectivity and diameter},
Internet Mathematics 11 (2015), 528--554.

\end{thebibliography}

\end{document}
